
\Needspace{12\baselineskip}
\subsubsection*{Naučnoistraživački rad -- \ClanDva}

\textbf{Naslov rada:}
\emph{Applying supervised machine learning algorithms and ensemble
models to enhance credit card fraud detection}.\\

\textbf{Autori:}
Abrar Al-Bulushi, Abdul K. Shaikh i Naresh Adhikari.\\

\textbf{Izvor:}
\emph{Frontiers in Artificial Intelligence}, vol. 9, članak 1813728,
2026. \cite{albulushi2026fraud}

\begin{table}[H]
\centering
\caption{Sažeti pregled naučnoistraživačkog rada Al-Bulushi i saradnika (2026)}
\label{tab:albulushi-pregled}
\small
\renewcommand{\arraystretch}{1.22}
\begin{tabularx}{\linewidth}{@{}p{0.25\linewidth}Y@{}}
\toprule
\textbf{Stavka} & \textbf{Opis} \\
\midrule

Problem &
Detekcija prevarnih transakcija kreditnim karticama kao binarna
klasifikacija na izrazito nebalansiranim podacima, uz provjeru
generalizacije modela na neviđenim skupovima podataka. \\

Skup podataka &
Glavni skup za treniranje i testiranje je \emph{isti Sparkov skup}
koji koristi naš tim (približno 1,3 miliona transakcija, 0,5\,\% prevara,
22 atributa). Skup od 555\,719 transakcija iz iste kolekcije i još četiri
javna skupa (među njima evropski skup od 284\,807 transakcija i PaySim)
služe kao neviđeni skupovi. \\

Korištene metode &
Logistička regresija, Naive Bayes, stablo odlučivanja, Random Forest,
neuronska mreža (ANN) i XGBoost; uravnotežavanje metodama RUS, ROS
i SMOTE; selekcija atributa (korelacija, ANOVA, RFE, unaprijedna
selekcija); novi atributi ponašanja korisnika; pomjeranje praga
odlučivanja; ansambli \emph{bagging}, \emph{boosting} i \emph{stacking}. \\

Metrike &
Tačnost, preciznost, osjetljivost (recall), F1-mjera i AUC,
uz konfuzijsku matricu; 4-struka unakrsna validacija. \\

Najbolji rezultat &
Pojedinačno: \textbf{XGBoost} bez uravnotežavanja, F1 \textbf{0,84}
(AUC 0,99). Konačni model \textbf{Bagging~1} (10 modela):
F1 \textbf{0,795}, recall 0,864, AUC 0,99; uz nove atribute ponašanja
F1 raste na \textbf{0,83} (recall 0,89). \\

\bottomrule
\end{tabularx}
\end{table}

\textbf{Cilj i istraživački problem.}
Autori razvijaju cjelovit okvir za detekciju prevara koji objedinjuje
izbor algoritma, uravnotežavanje klasa, selekciju i konstrukciju atributa
te ansambl modele. Posebno ih zanima da li model treniran na jednom
skupu zadržava performanse na drugim, neviđenim skupovima podataka.

\textbf{Skup podataka i priprema.}
Rad je za nas posebno relevantan jer kao glavni skup koristi upravo
Sparkov podatke sa Kaggle platforme, sa 12 kategoričkih i 11 numeričkih
atributa i bez nedostajućih vrijednosti. Podaci su standardizovani
(\emph{z-score}), a podjela je urađena 4-strukom unakrsnom validacijom
(75\,\% podataka za treniranje). Autori su konstruisali i nove atribute
ponašanja: oznaku da li je transakcija izvan uobičajenog vremenskog prozora
korisnika (izvan 10.--90.~percentila sati transakcija) i \emph{anomaly score}
dobijen metodom Isolation Forest.

\textbf{Poređenje eksperimentalnih rezultata.}
Tabela~\ref{tab:albulushi-rezultati} prikazuje izdvojene rezultate iz
Tabela 1, 8 i 9 rada (pojedinačni modeli, pomjeranje praga i ansambli).

\begin{table}[H]
\centering
\caption{Izdvojeni rezultati iz rada Al-Bulushi i saradnika (Sparkov skup)}
\label{tab:albulushi-rezultati}
\small
\setlength{\tabcolsep}{4pt}
\renewcommand{\arraystretch}{1.18}
\begin{tabularx}{\linewidth}{@{}Y c c c c c@{}}
\toprule
\textbf{Model} & \textbf{Acc.} & \textbf{F1} & \textbf{Prec.} & \textbf{Recall} & \textbf{AUC} \\
\midrule
Logistička regresija -- RUS & 0,87 & 0,06 & 0,03 & 0,76 & 0,96 \\
Naive Bayes -- bez uravnoteženja & 0,90 & 0,07 & 0,04 & 0,68 & 0,82 \\
Stablo odlučivanja -- ROS & 0,99 & 0,78 & 0,79 & 0,77 & 0,88 \\
Random Forest -- bez uravnoteženja & 0,99 & 0,82 & 0,88 & 0,76 & 0,97 \\
Random Forest -- RUS & 0,97 & 0,30 & 0,17 & 0,97 & 0,99 \\
ANN -- bez uravnoteženja & 0,99 & 0,83 & 0,93 & 0,76 & 0,99 \\
\textbf{XGBoost -- bez uravnoteženja} & \textbf{0,99} & \textbf{0,84} & \textbf{0,91} & \textbf{0,79} & \textbf{0,99} \\
XGBoost -- SMOTE, prag 0,7 & 0,99 & 0,70 & 0,67 & 0,85 & 0,99 \\
\midrule
Bagging~1 (10 modela) & 0,99 & 0,79 & 0,73 & 0,86 & 0,99 \\
Stacking~1 (RF baza, ANN meta) & 0,99 & 0,80 & 0,79 & 0,80 & 0,93 \\
Boosting~2 (LR + AdaBoost) & 0,99 & 0,00 & 0,00 & 0,00 & 0,82 \\
\bottomrule
\end{tabularx}
\end{table}

\textbf{Tumačenje rezultata.}
Logistička regresija i Naive Bayes daju vrlo nisku F1-mjeru (ispod 0,1)
bez obzira na metodu uravnotežavanja, dok modeli zasnovani na stablima
(Random Forest, XGBoost) i neuronska mreža postižu F1 oko 0,8.
Rezultati jasno pokazuju kompromis između preciznosti i osjetljivosti:
\emph{random undersampling} podiže recall Random Foresta na 0,97, ali
preciznost pada na 0,17, jer model treniran na uravnoteženim podacima
na realističnoj raspodjeli generiše mnogo lažnih alarma. Tačnost je
za gotovo sve modele 0,99 i zato nije informativna. Ansambl Bagging~1
nije postigao najveći F1, ali su ga autori izabrali kao konačni model zbog
stabilnosti i visokog recall-a, a na neviđenom Sparkov skupu postiže
F1 0,61, odnosno 0,66 uz nove atribute ponašanja.

\textbf{Zaključak i značaj za istraživanje.}
Rad potvrđuje da su na našem skupu podataka metode zasnovane na
stablima i \emph{boosting} najperspektivnije, a da su linearni modeli i
Naive Bayes slabi. Za naš projekat su posebno korisna tri zapažanja:
(1) uravnotežavanje treba primjenjivati samo na skupu za treniranje i
pratiti njegov uticaj na preciznost, (2) pomjeranje praga odlučivanja je
jednostavan način da se podesi odnos preciznosti i osjetljivosti, i
(3) atributi izvedeni iz vremena transakcije i ponašanja korisnika
mogu dodatno poboljšati detekciju.
